% Generated by Listas/make_statements.py from pratica-lecture-08.md, item 19. Do not edit.
\textbf{Source.} Mas-Colell, Whinston \& Green (1995), \emph{Microeconomic Theory}, Ch. 13 Exercises, Exercise 13.C.5\^{}B, p.~475 {[}PDF p.~248{]}. Reproduced verbatim. Answer: none in the book; entries for 13.C.1, 13.C.2, 13.C.4 and 13.C.5 are in the MWG solutions manual (\texttt{Livros/}).



\begin{quote}

\textbf{13.C.5\^{}B} Assume a single firm and a single consumer. The firm's product may be either high or low quality and is of high quality with probability \(\lambda\). The consumer cannot observe quality before purchase and is risk neutral. The consumer's valuation of a high-quality product is \(v_H\); her valuation of a low-quality product is \(v_L\). The costs of production for high (\(H\)) and low (\(L\)) quality are \(c_H\) and \(c_L\), respectively. The consumer desires at most one unit of the product. Finally, the firm's price is regulated and is set at \(p\). Assume that \(v_H > p > v_L > c_H > c_L\).



\begin{enumerate}

\def\labelenumi{(\alph{enumi})}

\item

  Given the level of \(p\), under what conditions will the consumer buy the product?

\item

  Suppose that before the consumer decides whether to buy, the firm (which knows its type) can advertise. Advertising conveys no information directly, but consumers can observe the total amount of money that the firm is spending on advertising, denoted by \(A\). Can there be a separating perfect Bayesian equilibrium, that is, an equilibrium in which the consumer rationally expects firms with different quality levels to pick different levels of advertising?

\end{enumerate}

\end{quote}

